% SPDX-FileCopyrightText: Copyright (c) 2016-2026 Objectionary.com
% SPDX-License-Identifier: MIT

\section{Worked Examples}\label{app:examples}

This appendix collects end-to-end walk-throughs of EO programs.
Each example exercises a different combination of the constructs introduced
  in \cref{sec:syntax,sec:wellformedness}, with reference to the
  \(\varphi\)-translation rules of \cref{sec:semantics}.

\subsection{Hello, world}\label{app:ex:hello}

The minimal greeting program decorates a copy of \ff{Q.io.stdout}.

\begin{ffcode}
+package org.eolang.examples

# Greets the world.
[] > app
  Q.io.stdout > @
    "Hello, world!\n"
\end{ffcode}

Construct usage, line by line:

\begin{itemize}
  \item Line 1 is a meta directive (\cref{sec:syntax:meta}) declaring the
    package. The file's top-level objects are children of
    \(\Phi.\mathtt{org}.\mathtt{eolang}.\mathtt{examples}\).
  \item Line 3 is a comment line (\cref{sec:syntax:comment}); it attaches
    to the named object below at the same indent.
  \item Line 4 is a top-level formation \ff{[] > app}: a master child with
    no voids. The blank line on line 2 is the required separator between
    the meta header and the program (\cref{sec:wf:blanks}).
  \item Line 5 is a vertical method chain \ff{Q.io.stdout} that becomes a
    vertical application as soon as its deeper-indent child appears.
    The suffix \ff{> @} names the expression as the decoratee of \ff{app}.
  \item Line 6 is a string literal argument of \ff{stdout}.
\end{itemize}

The \(\varphi\)-translation is
\begin{equation*}
  \Phi.\mathtt{org}.\mathtt{eolang}.\mathtt{examples}.\mathtt{app}\,\mapsto\,
    \llbracket\,
      \varphi \mapsto \Phi.\mathtt{io}.\mathtt{stdout}\,(\alpha_0 \mapsto \text{\ff{"Hello, world!\textbackslash n"}})
    \,\rrbracket.
\end{equation*}
Dataizing the program (\cref{sec:sem:atoms}) reduces the dispatch chain,
  invokes the \(\lambda\)-asset of \ff{stdout}, and writes the greeting
  bytes to the host's standard output.

\subsection{Mixed-construct Walkthrough}\label{app:ex:mixed}

The following program exercises eight different constructs in a single
  tree: meta header, comment block, formation, vertical method chain,
  star-tuple head, horizontal application, inline-phi formation, and
  bare head.

\begin{ffcode}
+architect someone@example.com
+version 1.0

# A program demonstrating most constructs.
# Comment block of two lines.
[args] > main
  mem.alloc > @
    0
    [x] >>
      run > @
        *
          x.put 2
          loop
            x.lt 6 > [i] >>
            [i] >>
              true
\end{ffcode}

Line-by-line classification (each line is annotated with the outer kind of
  the expression it introduces):

\begin{itemize}
  \item \ff{[args] > main} --- \ff{bare-formation},
    top-level master child of the package root.
    The comment block above attaches to this formation.
    By convention, the runtime applies a top-level \ff{main} to the
    program's command-line arguments, supplied as a tuple bound to the
    \ff{args} void; the convention is independent of the language and
    is honoured by the runtime, not by the parser.
  \item \ff{mem.alloc > @} --- \ff{hmethod} chain with 0 horizontal
    arguments and the name suffix \ff{> @}. As soon as the deeper-indent
    children appear, the outer kind transitions to \ff{vapplication}
    (\cref{sec:syntax:vapp}). The expression names the decoratee
    \(\varphi\) of \ff{main}.
  \item \ff{0} --- bare \ff{head}, supplying the first positional
    argument of \ff{mem.alloc}. No name suffix is required because the
    parent kind is \ff{hmethod}/\ff{vapplication}, not a formation.
  \item \ff{[x] >>} --- nested \ff{bare-formation} with an auto-generated
    name. The single void \(x\) is supplied by the caller.
  \item \ff{run > @} --- a \ff{head} with the name suffix \ff{> @};
    once the deeper block opens, it becomes a \ff{vapplication}.
  \item \ff{*} --- a star tuple as the head of an application
    (\cref{sec:syntax:app}); with deeper children it becomes a
    \(\Phi.\mathtt{tuple}\) constructor over those children.
  \item \ff{x.put 2} --- \ff{happlication}, a horizontal application
    whose head is the hmethod \ff{x.put}. Horizontally-completed
    (\cref{sec:wf:notation}): no deeper children, no \ff{.method} wrap.
  \item \ff{loop} --- a bare \ff{head}; opens a child block as a
    \ff{vapplication}.
  \item \ff{x.lt 6 > [i] >>} --- an \ff{inline-phi-formation}.
    The left-hand expression \ff{x.lt 6} becomes the \(\varphi\)-bound
    body of a fresh formation with one void \(i\); the formation receives
    an auto-generated name.
  \item \ff{[i] >>} --- another \ff{bare-formation} with one void and an
    auto-name.
  \item \ff{true} --- a bare \ff{head}, the body of the inner formation.
\end{itemize}

The example illustrates how name requirements vary by parent kind:
  the child \ff{0} carries no suffix because its parent is a vertical
  application (which positional-indexes its arguments), while the child
  \ff{x.lt 6 > [i] >>} must be named because the inline-phi composite
  always introduces a named object (\cref{sec:wf:naming}).

\subsection{Leap-year predicate}\label{app:ex:leap}

This example exercises reversed dispatch in both horizontal and vertical
  form, paren groups as receivers, the const-marker, and the inline
  binding \ff{:y}.

\begin{ffcode}
[args] > main
  [y] > leap
    or. > @
      and.
        eq. (mod. y 4) 0
        not. (eq. (mod. y 100) 0)
      eq. (mod. y 400) 0
  out. > @
    format
      "%d is %sa leap year!"
      (args.at 0).as-int > year!
      if (leap year:y) "" "not "
\end{ffcode}

Construct highlights:

\begin{itemize}
  \item \ff{[y] > leap} --- a nested formation with one void; the body
    computes the leap-year condition.
  \item \ff{or. > @} --- a vertical \ff{bare-reversed} dispatch.
    The first deeper child is the receiver of \ff{.or}; subsequent
    children are method arguments.
  \item \ff{and.} --- another vertical \ff{bare-reversed}, used as the
    receiver of \ff{or.}.
  \item \ff{eq. (mod. y 4) 0} --- a horizontal \ff{reversed-with-hargs}.
    The receiver is the paren group \ff{(mod. y 4)}; the method argument
    is \ff{0}. The line is horizontally-completed
    (\cref{sec:wf:notation}).
  \item \ff{not. (eq. (mod. y 100) 0)} --- a horizontal reversed dispatch
    with a single argument (the receiver); the receiver is itself a
    horizontal reversed dispatch.
  \item \ff{out. > @} --- another vertical reversed dispatch on the
    top-level formation, naming the decoratee of \ff{main}.
  \item \ff{(args.at 0).as-int > year!} --- an \ff{hmethod} chain whose
    head is a paren group, with the const-marked name \ff{year!}.
    The intermediate vapplication argument is named, demonstrating that
    vertical-application arguments may each carry their own optional names.
  \item \ff{if (leap year:y) "" "not "} --- an \ff{happlication} with
    three horizontal arguments. The first argument is the paren group
    \ff{(leap year:y)}, an application of \ff{leap} to \ff{year} with
    the inline binding \ff{:y}; the remaining arguments are string
    literals.
\end{itemize}

The receivers of reversed dispatches are exempt from the all-or-nothing
  binding rule (\cref{sec:wf:bindings}); only the method arguments form
  the binding group.

\subsection{Multi-line Bytes Literal}\label{app:ex:bytes}

A \ff{BYTES} literal of four bytes spread across two source lines:

\begin{ffcode}
size.
  CA-FE-
  BE-BE
\end{ffcode}

The trailing dash on the first chunk continues the literal onto the next
  line.
The two indented lines under \ff{size.} form a single bytes token
  \ff{CA-FE-BE-BE} occupying one argument slot of the reversed dispatch
  \ff{size.}.
The continuation is a lexical concern (\cref{sec:lexical:bytes}); the
  indent stack sees one expression at one indent.

\subsection{Compact Tuple}\label{app:ex:tuple}

The compact-tuple line \ff{seq * > @} folds its three vertical children
  into a single \(\Phi.\mathtt{tuple}\) argument of \ff{seq}:

\begin{ffcode}
seq * > @
  step-1
  step-2
  step-3
\end{ffcode}

translates to
\begin{equation*}
  \varphi\,\mapsto\,
    \mathtt{seq}\,(
      \alpha_0 \mapsto \Phi.\mathtt{tuple}(\mathtt{step\text{-}1},\ \mathtt{step\text{-}2},\ \mathtt{step\text{-}3})
    ).
\end{equation*}

A non-zero \(N\) splits the children between direct positional arguments
  and the tuple wrapper:

\begin{ffcode}
sprintf *1 > greeting
  "Hello, %s!"
  "Jeff"
  "extra"
\end{ffcode}

binds \ff{"Hello, \%s!"} as \(\alpha_0\) and folds the remaining two
  strings into \(\alpha_1 \mapsto \Phi.\mathtt{tuple}(\mathtt{"Jeff"},\,\mathtt{"extra"})\).

\subsection{Wrapping a Vertical-completed Expression}\label{app:ex:vmethod-wrap}

A same-indent \ff{.method} line may wrap a multi-line expression whose
  body block has already closed.
The wrap transitions the underlying expression from
  \emph{vertical-completed} back to \emph{open} as a \ff{vmethod}
  (\cref{sec:wf:notation}).

\begin{ffcode}
[] > foo
  bar
    a
    b
  .baz > result
\end{ffcode}

The head line \ff{bar} opens a \ff{vapplication} whose body is the
  block \ff{a, b}.
The deeper-indent block closes when the next line at indent 2 arrives;
  that line is \ff{.baz > result}.
The dispatch on \ff{.baz} wraps the just-closed \ff{vapplication} into
  a \ff{vmethod}, naming the chain \ff{result}.
The \(\varphi\)-translation is
\begin{equation*}
  \mathtt{result}\;\mapsto\;
    \mathtt{bar}\,(\alpha_0 \mapsto \mathtt{a},\,\alpha_1 \mapsto \mathtt{b})\,.\,\mathtt{baz}.
\end{equation*}

A second same-indent \ff{.method} would extend the chain further.
A \emph{horizontally-completed} expression cannot be wrapped this way:
  the line \ff{bar a b} followed by a same-indent \ff{.baz} is rejected
  because \ff{happlication} is in the set \(\mathcal{H}\)
  (\cref{sec:wf:notation}).

\subsection{Inline Binding on Non-final Chain Link}\label{app:ex:binding-chain}

An inline binding \ff{:label} or \ff{:N} on a method dispatch is legal
  only when the method is the \emph{last} link in its chain
  (\cref{sec:wf:bindings}).
A binding on an intermediate link is rejected.

\begin{ffcode}
foo (bar.baz:0 1)        // rejected: :0 attaches to bar.baz,
                         //   but .baz is not the chain's last link
foo (bar.baz 1):0        // legal: :0 attaches to the whole bar.baz 1
foo (bar.baz.qux:0 1)    // rejected: .baz is mid-chain, :0 forbidden
foo (bar.baz.qux 1:0)    // legal: :0 attaches to the harg 1 of .qux,
                         //   the last link
\end{ffcode}

The rule prevents an ambiguity in argument attachment: a binding on a
  mid-chain link could plausibly refer to that link's receiver or to
  the chain's tail.
Restricting the binding to the last link makes the attachment
  unambiguous.

\subsection{Nested Atom with Empty Body}\label{app:ex:nested-atom}

An atom that does not sit at indent 0 is a \emph{nested atom}.
Test attributes are admissible only at indent level 1 of a top-level
  object (\cref{sec:wf:atoms}), so a nested atom cannot hold tests;
  an atom admits only test children, so a nested atom cannot hold
  plain or master children either.
The body of a nested atom is therefore \emph{empty}.

\begin{ffcode}
[] > outer
  [] > middle
    [a b] > inner /number   // legal nested atom; body must be empty
\end{ffcode}

The line \ff{[a b] > inner /number} declares a formation with two
  voids and the atom signature \ff{/number}.
Its body is absent; the formation closes at end-of-block.
The \(\varphi\)-translation binds \ff{inner} inside \ff{middle} with
  the two voids and a \(\lambda\)-asset whose host-language function
  returns a value of type \ff{number}.

A nested atom inside another atom is rejected outright:

\begin{ffcode}
[] > outer /number
  [] > inner /integer       // rejected: atom inside an atom
\end{ffcode}

The outer atom admits only test attributes; an atom is a master child,
  not a test attribute, and is therefore inadmissible.

\subsection{Dangling Comment}\label{app:ex:dangling-comment}

A comment block must be immediately followed by a named object at the
  same indent (\cref{sec:wf:comments}).
A block followed by a blank line, by a deeper- or shallower-indent
  object, or by no named object at all, is \emph{dangling} and is
  rejected.

\begin{ffcode}
# This comment leads nowhere.

[] > app                     // rejected: blank line breaks attachment;
                             //   the comment above is dangling.
\end{ffcode}

A blank line between the comment block and the following object breaks
  attachment regardless of indent.
Two comment blocks separated by a blank line at the same indent are
  not merged: the earlier block is dangling, and only the later block
  attaches.

\begin{ffcode}
# First block.

# Second block.
[] > app                     // legal: second block attaches; first dangling
\end{ffcode}

A comment block at indent \(K\) followed by a named object at a
  different indent is also dangling:

\begin{ffcode}
# Attached to whom?
  [] > inner                 // rejected: comment at indent 0 cannot
                             //   attach to an object at indent 2
\end{ffcode}
